\section{Quantitative Study}
\label{sec:experiments}

\subsection{Experimental Setup}

We evaluate on all $2{,}997$ Img2GPS3k images~\cite{vo2017revisiting}.
WeDetect-Uni Base generates category-agnostic proposals once for the complete
set. We cache boxes at an objectness threshold of $0.15$ after non-maximum
suppression ($0.7$ IoU), then form the intervention mask from the union of at
most $100$ boxes scoring at least $0.4$. This produces $6{,}152$ boxes in
$2{,}081$ images; $2{,}058$ images retain a non-empty mask after CLIP's centre
crop. Images without an active mask remain in the evaluation as exact no-ops,
so the reported metrics describe the complete benchmark.

For every active layer, the same in-region and out-of-region biases are applied
to all $16$ heads, with base magnitudes $a=+2$ and $b=-2$. We measure top-1
geodesic error and
prediction displacement. For image $i$, improvement is
$\Delta_i=d_i^{\mathrm{base}}-d_i^{\mathrm{int}}$, so positive values are
better. Because a gallery change can move a prediction across continents, the
selection score is the mean of
$\operatorname{clip}(\Delta_i,-2500,2500)$~km. Confidence intervals use
$10{,}000$ paired bootstrap resamples over images.

We deterministically assign $1{,}079$ images to discovery and $1{,}918$ to
holdout by hashing the file name. Model selection uses discovery only. A
preliminary mixed-precision layer sweep on the discovery split supplies five candidates,
$\{2,4,7,12,3\}$, ordered by discovery score. We evaluate all $26$ subsets containing two to five of
these layers under three strength schedules: fixed $2$ per layer,
$2/\sqrt{k}$, and $2/k$, where $k$ is the number of active layers. Together
with five single-layer controls, this gives $78$ combinations and five
controls. The confirmatory pass uses deterministic FP32 with TF32 disabled.

\subsection{Held-Out Results}

The discovery split selects layers $\{2,7,12\}$ with the $2/k$ schedule,
equivalent to magnitude $0.667$ at each layer. Table~\ref{tab:main} compares
this configuration with the discovery-selected single-layer control. The
apparent discovery gain does not transfer: the selected combination changes
from $+11.94$~km on discovery to $-3.66$~km on holdout. Its paired advantage
over layer 2 is only $0.64$~km (95\% CI $[-16.83,17.57]$). The experiment
therefore does not support synergy between individually promising layers.

\begin{table}[t]
\centering
\caption{Deterministic FP32 results. $\Delta_c$ is mean improvement clipped to
$\pm2500$~km; positive is better. Columns after the discovery score are holdout
metrics. Every column but the last is in kilometres.}
\label{tab:main}
\small
\setlength{\tabcolsep}{2.5pt}
\begin{tabular}{l|rrrrrr}
\toprule
\textbf{Configuration} &
\textbf{Disc. $\Delta_c$} & \textbf{Error} & \textbf{Hold. $\Delta_c$} &
\textbf{95\% CI} & \textbf{Shift} & \textbf{Cosine} \\
\midrule
GeoCLIP baseline & 0.00 & 1716.31 & 0.00 & --- & 0.00 & 1.0000 \\
Layer 2 & +4.95 & 1731.44 & -4.30 & $[-18.83,9.96]$ & 136.15 & 0.9993 \\
Layers 2, 7, 12 ($2/k$) & +11.94 & 1728.86 & -3.66 & $[-20.43,12.40]$ & 189.53 & 0.9984 \\
\bottomrule
\end{tabular}
\end{table}

Despite the absence of an accuracy gain, the intervention exposes output
sensitivity. The selected configuration preserves a mean cosine similarity of
$0.9984$ to the baseline image embedding, yet moves the predicted coordinate
by $189.5$~km on average. GeoCLIP retrieves from $10^5$ location embeddings;
small movements in representation space can therefore cross a gallery
decision boundary and produce a geographically large change.

Increasing the number of layers does not help on average. With fixed magnitude
$2$, mean holdout $\Delta_c$ declines from $-8.22$~km for two layers to
$-39.35$~km for five. The $2/k$ schedule limits this accumulation ($-2.01$ to
$-0.13$~km across sizes) without yielding a consistent positive effect. Candidate rankings also transfer
weakly from discovery to holdout ($r=0.259$), underscoring the risk of
reporting the best result on the same images used to find it.

This conclusion is not an artefact of diluting the result with no-op images.
On the $1{,}321$ holdout images with an active mask, the selected configuration
obtains $-5.32$~km (95\% CI $[-28.75,18.02]$), and its paired advantage over
layer 2 remains only $0.92$~km (95\% CI $[-23.66,25.66]$).

A natural objection is that displacement merely tracks how much of the image
the mask covers. It does not. Over the $2{,}051$ active-mask images of the
all-layer $a=2$, $b=-2$ run, coverage and displacement are negatively related
(Spearman $\rho=-0.256$): the least-covered quartile moves the prediction by a
median of $407$~km with $13.8\%$ of images unmoved, whereas the quartile
covering more than $81\%$ of patches moves it by a median of $0$~km with
$49.4\%$ unmoved. Eq.~\ref{eq:bias} predicts this: a mask spanning nearly
every patch biases every patch key almost equally, and softmax is invariant to
a constant shift. The experiment still isolates
layer-combination effects rather than showing that WeDetect regions are more
informative than arbitrary ones; that claim needs an area-matched random-mask
control.
